% lang/pl.tex
% -----------------------------------------------------------------------
% "Poznaję ekonomię": las cadenas de texto del cuaderno en polaco, las
% mismas que lang/es.tex, una a una. polish.tex fija \booklang{polish}
% antes de % preamble.tex
% -----------------------------------------------------------------------
% "Aprendo economía" -- cuaderno diario de economía (español, A4). Todo
% el aparato tipográfico vive aquí: paquetes, colores, cajas y las macros
% que usa cada página. Es el motor de "Aprendo a leer"
% (https://github.com/konradcinkusz/learning-to-read) y de "Aprendo los
% números" (https://github.com/konradcinkusz/learning-to-count) -- la
% misma letra, la misma paleta, la misma cabecera de cada día, las mismas
% cajas --, con lo que cambia: la caja de arriba es "Hoy" (la historia
% del día y un dibujo, y los lunes la palabra de la semana), y las
% actividades son las de la economía de cada día (ver notes/01-plan.md).
%
% Requiere LuaLaTeX (fontspec): la letra es Andika (SIL, OFL,
% fonts/andika/), pensada para primeros lectores.
% -----------------------------------------------------------------------

\usepackage{fontspec}
\setmainfont{Andika}[
  Path = fonts/andika/,
  Extension = .ttf,
  UprightFont = *-Regular,
  BoldFont = *-Bold,
  ItalicFont = *-Italic,
  BoldItalicFont = *-BoldItalic,
]
\setsansfont{Andika}[
  Path = fonts/andika/,
  Extension = .ttf,
  UprightFont = *-Regular,
  BoldFont = *-Bold,
  ItalicFont = *-Italic,
  BoldItalicFont = *-BoldItalic,
]

\usepackage[spanish]{babel}
% babel-spanish hace de " un carácter activo; las comillas escritas a mano
% en el texto se estropean sin esto (ver el mismo comentario, más largo,
% en el preamble.tex de "Aprendo a leer"). Solo funciona así, al empezar
% el documento.
\AtBeginDocument{\shorthandoff{"}}

% El pie de página ("Día N / 260") vive en el pie de verdad, no en la
% página de texto: la caja de actividad llena todo lo que queda de página
% (ver `cajadia`), y un pie dentro del texto la empujaría a otra página.
\usepackage[a4paper,top=18mm,bottom=22mm,left=22mm,right=18mm,footskip=10mm]{geometry}

\usepackage{xcolor}
\usepackage[most]{tcolorbox}
\usepackage{tabularx}
\usepackage{booktabs}
\usepackage{enumitem}
\usepackage{tikz}
\usetikzlibrary{arrows.meta}
\usepackage{graphicx}
\usepackage{multicol}

\pagestyle{empty}
\setlength{\parindent}{0pt}
\setlength{\parskip}{2pt}

% --- cadenas de texto ----------------------------------------------------
% lang/es.tex
% -----------------------------------------------------------------------
% Todas las cadenas de texto del cuaderno, en un solo sitio, como en
% "Aprendo a leer": cambiar una palabra aquí la cambia en las 260 páginas
% a la vez, sin tocar el contenido generado.
% -----------------------------------------------------------------------

% Total de días del cuaderno. tools/gen_economia.py --check comprueba que
% coincide con su propio TOTAL_DIAS.
\newcommand{\totaldias}{260}

% La cabecera de cada día (\cabeceraDia, preamble.tex).
\newcommand{\lblFecha}{Fecha}
\newcommand{\lblNotas}{Notas}
\newcommand{\lblDia}{Día}
\newcommand{\lblSemana}{Semana}
\newcommand{\lblTrimestre}{Trimestre}
\newcommand{\lblHecho}{Hecho}
\newcommand{\lblHechoSola}{sola}
\newcommand{\lblHechoConAyuda}{con ayuda}
\newcommand{\lblHechoConDificultad}{con dificultad}

% La caja de arriba, la misma todos los días.
\newcommand{\lblHoy}{Hoy}
\newcommand{\lblPalabraSemana}{La palabra de la semana}

% Los títulos de las cajas de actividad.
\newcommand{\lblClasifica}{Clasifica}
\newcommand{\lblUne}{Une}
\newcommand{\lblVF}{¿Verdad o mentira?}
\newcommand{\lblTrueque}{Trueque}
\newcommand{\lblDibuja}{Dibuja}
\newcommand{\lblRepasa}{Repasa}

% Las dos letras de "¿Verdad o mentira?".
\newcommand{\lblVerdad}{V}
\newcommand{\lblMentira}{F}

% La portada.
\newcommand{\lblTituloPortada}{Aprendo economía}
\newcommand{\lblSubtituloPortada}{El dinero, el trabajo y las decisiones de cada día}
\newcommand{\lblSubtituloPortadaMetodo}{Cuaderno diario para aprender a ahorrar, elegir, cambiar y compartir -- con la familia de \emph{Aprendo a leer}}
\newcommand{\lblPortadaNombre}{Nombre}
\newcommand{\lblPortadaCurso}{Curso}

% La medalla del final de cada trimestre (T1-T3) y el cartel de cada
% diez páginas.
\newcommand{\lblMedalla}[1]{¡Medalla de #1!}
\newcommand{\lblMedallaAnimo}{¡Sigue así, \rule{55mm}{0.4pt}!}
\newcommand{\lblPaginasHechas}[1]{¡#1 páginas hechas!}

% La clave de respuestas y el diccionario.
\newcommand{\lblClave}{Clave de respuestas}
\newcommand{\lblDiccionario}{Mi diccionario de economía}


% =========================================================================
% Paleta
% =========================================================================
% La de "Aprendo a leer": cada actividad tiene un color saturado (borde y
% título) y uno claro (fondo), y el color es un canal EXTRA, nunca el
% único -- el título en negrita ya dice qué actividad es. main-bw.tex fija
% \bookcolor{bw} antes de % preamble.tex
% -----------------------------------------------------------------------
% "Aprendo economía" -- cuaderno diario de economía (español, A4). Todo
% el aparato tipográfico vive aquí: paquetes, colores, cajas y las macros
% que usa cada página. Es el motor de "Aprendo a leer"
% (https://github.com/konradcinkusz/learning-to-read) y de "Aprendo los
% números" (https://github.com/konradcinkusz/learning-to-count) -- la
% misma letra, la misma paleta, la misma cabecera de cada día, las mismas
% cajas --, con lo que cambia: la caja de arriba es "Hoy" (la historia
% del día y un dibujo, y los lunes la palabra de la semana), y las
% actividades son las de la economía de cada día (ver notes/01-plan.md).
%
% Requiere LuaLaTeX (fontspec): la letra es Andika (SIL, OFL,
% fonts/andika/), pensada para primeros lectores.
% -----------------------------------------------------------------------

\usepackage{fontspec}
\setmainfont{Andika}[
  Path = fonts/andika/,
  Extension = .ttf,
  UprightFont = *-Regular,
  BoldFont = *-Bold,
  ItalicFont = *-Italic,
  BoldItalicFont = *-BoldItalic,
]
\setsansfont{Andika}[
  Path = fonts/andika/,
  Extension = .ttf,
  UprightFont = *-Regular,
  BoldFont = *-Bold,
  ItalicFont = *-Italic,
  BoldItalicFont = *-BoldItalic,
]

\usepackage[spanish]{babel}
% babel-spanish hace de " un carácter activo; las comillas escritas a mano
% en el texto se estropean sin esto (ver el mismo comentario, más largo,
% en el preamble.tex de "Aprendo a leer"). Solo funciona así, al empezar
% el documento.
\AtBeginDocument{\shorthandoff{"}}

% El pie de página ("Día N / 260") vive en el pie de verdad, no en la
% página de texto: la caja de actividad llena todo lo que queda de página
% (ver `cajadia`), y un pie dentro del texto la empujaría a otra página.
\usepackage[a4paper,top=18mm,bottom=22mm,left=22mm,right=18mm,footskip=10mm]{geometry}

\usepackage{xcolor}
\usepackage[most]{tcolorbox}
\usepackage{tabularx}
\usepackage{booktabs}
\usepackage{enumitem}
\usepackage{tikz}
\usetikzlibrary{arrows.meta}
\usepackage{graphicx}
\usepackage{multicol}

\pagestyle{empty}
\setlength{\parindent}{0pt}
\setlength{\parskip}{2pt}

% --- cadenas de texto ----------------------------------------------------
% lang/es.tex
% -----------------------------------------------------------------------
% Todas las cadenas de texto del cuaderno, en un solo sitio, como en
% "Aprendo a leer": cambiar una palabra aquí la cambia en las 260 páginas
% a la vez, sin tocar el contenido generado.
% -----------------------------------------------------------------------

% Total de días del cuaderno. tools/gen_economia.py --check comprueba que
% coincide con su propio TOTAL_DIAS.
\newcommand{\totaldias}{260}

% La cabecera de cada día (\cabeceraDia, preamble.tex).
\newcommand{\lblFecha}{Fecha}
\newcommand{\lblNotas}{Notas}
\newcommand{\lblDia}{Día}
\newcommand{\lblSemana}{Semana}
\newcommand{\lblTrimestre}{Trimestre}
\newcommand{\lblHecho}{Hecho}
\newcommand{\lblHechoSola}{sola}
\newcommand{\lblHechoConAyuda}{con ayuda}
\newcommand{\lblHechoConDificultad}{con dificultad}

% La caja de arriba, la misma todos los días.
\newcommand{\lblHoy}{Hoy}
\newcommand{\lblPalabraSemana}{La palabra de la semana}

% Los títulos de las cajas de actividad.
\newcommand{\lblClasifica}{Clasifica}
\newcommand{\lblUne}{Une}
\newcommand{\lblVF}{¿Verdad o mentira?}
\newcommand{\lblTrueque}{Trueque}
\newcommand{\lblDibuja}{Dibuja}
\newcommand{\lblRepasa}{Repasa}

% Las dos letras de "¿Verdad o mentira?".
\newcommand{\lblVerdad}{V}
\newcommand{\lblMentira}{F}

% La portada.
\newcommand{\lblTituloPortada}{Aprendo economía}
\newcommand{\lblSubtituloPortada}{El dinero, el trabajo y las decisiones de cada día}
\newcommand{\lblSubtituloPortadaMetodo}{Cuaderno diario para aprender a ahorrar, elegir, cambiar y compartir -- con la familia de \emph{Aprendo a leer}}
\newcommand{\lblPortadaNombre}{Nombre}
\newcommand{\lblPortadaCurso}{Curso}

% La medalla del final de cada trimestre (T1-T3) y el cartel de cada
% diez páginas.
\newcommand{\lblMedalla}[1]{¡Medalla de #1!}
\newcommand{\lblMedallaAnimo}{¡Sigue así, \rule{55mm}{0.4pt}!}
\newcommand{\lblPaginasHechas}[1]{¡#1 páginas hechas!}

% La clave de respuestas y el diccionario.
\newcommand{\lblClave}{Clave de respuestas}
\newcommand{\lblDiccionario}{Mi diccionario de economía}


% =========================================================================
% Paleta
% =========================================================================
% La de "Aprendo a leer": cada actividad tiene un color saturado (borde y
% título) y uno claro (fondo), y el color es un canal EXTRA, nunca el
% único -- el título en negrita ya dice qué actividad es. main-bw.tex fija
% \bookcolor{bw} antes de % preamble.tex
% -----------------------------------------------------------------------
% "Aprendo economía" -- cuaderno diario de economía (español, A4). Todo
% el aparato tipográfico vive aquí: paquetes, colores, cajas y las macros
% que usa cada página. Es el motor de "Aprendo a leer"
% (https://github.com/konradcinkusz/learning-to-read) y de "Aprendo los
% números" (https://github.com/konradcinkusz/learning-to-count) -- la
% misma letra, la misma paleta, la misma cabecera de cada día, las mismas
% cajas --, con lo que cambia: la caja de arriba es "Hoy" (la historia
% del día y un dibujo, y los lunes la palabra de la semana), y las
% actividades son las de la economía de cada día (ver notes/01-plan.md).
%
% Requiere LuaLaTeX (fontspec): la letra es Andika (SIL, OFL,
% fonts/andika/), pensada para primeros lectores.
% -----------------------------------------------------------------------

\usepackage{fontspec}
\setmainfont{Andika}[
  Path = fonts/andika/,
  Extension = .ttf,
  UprightFont = *-Regular,
  BoldFont = *-Bold,
  ItalicFont = *-Italic,
  BoldItalicFont = *-BoldItalic,
]
\setsansfont{Andika}[
  Path = fonts/andika/,
  Extension = .ttf,
  UprightFont = *-Regular,
  BoldFont = *-Bold,
  ItalicFont = *-Italic,
  BoldItalicFont = *-BoldItalic,
]

\usepackage[spanish]{babel}
% babel-spanish hace de " un carácter activo; las comillas escritas a mano
% en el texto se estropean sin esto (ver el mismo comentario, más largo,
% en el preamble.tex de "Aprendo a leer"). Solo funciona así, al empezar
% el documento.
\AtBeginDocument{\shorthandoff{"}}

% El pie de página ("Día N / 260") vive en el pie de verdad, no en la
% página de texto: la caja de actividad llena todo lo que queda de página
% (ver `cajadia`), y un pie dentro del texto la empujaría a otra página.
\usepackage[a4paper,top=18mm,bottom=22mm,left=22mm,right=18mm,footskip=10mm]{geometry}

\usepackage{xcolor}
\usepackage[most]{tcolorbox}
\usepackage{tabularx}
\usepackage{booktabs}
\usepackage{enumitem}
\usepackage{tikz}
\usetikzlibrary{arrows.meta}
\usepackage{graphicx}
\usepackage{multicol}

\pagestyle{empty}
\setlength{\parindent}{0pt}
\setlength{\parskip}{2pt}

% --- cadenas de texto ----------------------------------------------------
\input{lang/es}

% =========================================================================
% Paleta
% =========================================================================
% La de "Aprendo a leer": cada actividad tiene un color saturado (borde y
% título) y uno claro (fondo), y el color es un canal EXTRA, nunca el
% único -- el título en negrita ya dice qué actividad es. main-bw.tex fija
% \bookcolor{bw} antes de \input{preamble}, y los mismos nombres de color
% apuntan a negro y gris: la misma página, sin tinta de color.
\makeatletter
\newif\ifbwbook
\def\aln@bookcolor@bw{bw}
\@ifundefined{bookcolor}{\bwbookfalse}{%
  \ifx\bookcolor\aln@bookcolor@bw \bwbooktrue \else \bwbookfalse \fi
}
\makeatother

\ifbwbook
  \definecolor{colorLectura}{gray}{0}
  \definecolor{colorLecturaClaro}{gray}{0.95}
  \definecolor{colorDibuja}{gray}{0}
  \definecolor{colorDibujaClaro}{gray}{0.95}
  \definecolor{colorCompleta}{gray}{0}
  \definecolor{colorCompletaClaro}{gray}{0.95}
  \definecolor{colorResponde}{gray}{0}
  \definecolor{colorRespondeClaro}{gray}{0.95}
  \definecolor{colorRelaciona}{gray}{0}
  \definecolor{colorRelacionaClaro}{gray}{0.95}
  \definecolor{colorCrea}{gray}{0}
  \definecolor{colorCreaClaro}{gray}{0.95}
  \definecolor{colorGris}{gray}{0.38}
  \definecolor{colorPunto}{gray}{0.25}
\else
  \definecolor{colorLectura}{HTML}{2E7D32}
  \definecolor{colorLecturaClaro}{HTML}{EAF6EC}
  \definecolor{colorDibuja}{HTML}{1565C0}
  \definecolor{colorDibujaClaro}{HTML}{E7F0FB}
  \definecolor{colorCompleta}{HTML}{C2185B}
  \definecolor{colorCompletaClaro}{HTML}{FBE9EF}
  \definecolor{colorResponde}{HTML}{6A1B9A}
  \definecolor{colorRespondeClaro}{HTML}{F2E9F7}
  \definecolor{colorRelaciona}{HTML}{E65100}
  \definecolor{colorRelacionaClaro}{HTML}{FDEEE2}
  \definecolor{colorCrea}{HTML}{AD1457}
  \definecolor{colorCreaClaro}{HTML}{FBE7EE}
  \definecolor{colorGris}{HTML}{616161}
  \definecolor{colorPunto}{HTML}{6A1B9A}
\fi

% =========================================================================
% Los dibujos
% =========================================================================
% diagrams/kit.tex: las piezas de los dibujos de «First Words» (de
% "Aprendo a leer"); diagrams/objetos.tex: las cosas de "Aprendo los
% números", cada una en la misma caja, con las monedas y los billetes;
% diagrams/economia.tex: las cosas nuevas de este cuaderno (la lupa, los
% cromos, las canicas, la hucha...).
\input{diagrams/kit}
\input{diagrams/objetos}
\input{diagrams/economia}

% =========================================================================
% Cajas
% =========================================================================
% La caja de actividad llena todo lo que queda de página (`height fill`),
% como en los otros cuadernos. Lo que tiene que ir abajo del todo (la
% línea de escribir, el cartel de cada diez páginas) va en la parte de
% abajo de la caja (\tcblower, estilo `abajo`); lo que tiene que ir
% centrado en el hueco, con `centrado abajo`.
\tcbset{
  cajabase/.style={
    enhanced,
    boxrule=1.1pt,
    arc=3mm,
    left=6mm, right=6mm, top=5mm, bottom=5mm,
    fonttitle=\bfseries\sffamily\large,
    coltitle=white,
  },
  cajadia/.style={cajabase, height fill},
  abajo/.style={space to upper, segmentation hidden, middle=2mm},
  centrado abajo/.style={space to lower, valign lower=center, segmentation hidden, middle=2mm},
}

\newtcolorbox{cajaHoy}{cajabase, top=3mm, bottom=3mm,
  colback=colorLecturaClaro, colframe=colorLectura, colbacktitle=colorLectura,
  title=\lblHoy}

% Una caja de actividad: #1 el nombre (cajaClasifica...), #2 el título
% (\lblClasifica...), #3 el color.
\newcommand{\cajaActividad}[3]{%
  \newtcolorbox{#1}[1][]{cajadia, ##1,
    colback=#3Claro, colframe=#3, colbacktitle=#3, title=#2}}
\cajaActividad{cajaClasifica}{\lblClasifica}{colorRelaciona}
\cajaActividad{cajaUne}{\lblUne}{colorRelaciona}
\cajaActividad{cajaVF}{\lblVF}{colorResponde}
\cajaActividad{cajaTrueque}{\lblTrueque}{colorCompleta}
\cajaActividad{cajaDibuja}{\lblDibuja}{colorDibuja}
\cajaActividad{cajaRepasa}{\lblRepasa}{colorCrea}

% La instrucción de cada actividad: la lee el adulto.
\newcommand{\instruccion}[1]{{\footnotesize\color{colorGris}#1\par}}
% El enunciado grande, lo que se hace: "Une cada cosa con su caja." Sin
% justificar y sin partir palabras, como la historia de "Hoy": lo lee la
% niña o el niño.
\newcommand{\sinCortes}{\hyphenpenalty=10000\exhyphenpenalty=10000}
\newcommand{\sinPartir}{\raggedright\sinCortes}
\newcommand{\enunciado}[1]{{\LARGE\sinPartir #1\par}}
% El hueco en el que se escribe la respuesta (#1, el ancho).
\newcommand{\huecoRespuesta}[1][26mm]{\raisebox{-5mm}{\tikz\draw[line width=1.1pt,
  rounded corners=2mm, fill=white] (0,0) rectangle (#1,18mm);}}
% Lo que no cabe en su sitio, encogido hasta que quepa (#1, el ancho).
\newsavebox{\cajaAjustada}
\newcommand{\ajustado}[2]{\sbox{\cajaAjustada}{#2}%
  \ifdim\wd\cajaAjustada>#1\relax\resizebox{#1}{!}{\usebox{\cajaAjustada}}%
  \else\usebox{\cajaAjustada}\fi}
% El cartel de cada diez páginas, al pie de "Repasa".
\newcommand{\cartel}[1]{\begin{center}\bfseries\Large\color{colorCrea}#1\end{center}}

\newcommand{\casilla}{\raisebox{-0.15ex}{\framebox[3.4mm][c]{\rule{0pt}{2.6mm}}}}
\newenvironment{listaRepaso}{%
  \begin{itemize}[label=\casilla, leftmargin=8mm, itemsep=2mm, topsep=1mm]\large
}{%
  \end{itemize}
}

% "¿Verdad o mentira?": cada frase, con su V y su F para rodear.
\newcommand{\letraVF}[1]{\tikz[baseline=-1.1mm]{\draw[line width=1pt] (0,0) circle (4.2mm);
  \node[font=\Large\bfseries] at (0,0) {#1};}}
\newcommand{\frasevf}[2]{%
  \noindent\begin{tabularx}{\linewidth}{@{}r@{\hspace{3mm}}>{\sinPartir\arraybackslash\Large}X@{\hspace{5mm}}c@{\hspace{3mm}}c@{}}
  {\Large\bfseries #1.} & #2 & \letraVF{\lblVerdad} & \letraVF{\lblMentira} \\
  \end{tabularx}\par\vspace{9mm}}

% =========================================================================
% La página de un día
% =========================================================================
% La cabecera es la de "Aprendo a leer": el día, la semana y el
% trimestre; el tema de la semana (el mismo que esa semana en "Leo con
% lupa"), y los campos para el adulto. Deja \label{dia:N}:
% tools/check_pages.py comprueba con esas etiquetas que cada día ocupa
% exactamente una página.
\newcommand{\cabeceraDia}[4]{%
  \label{dia:#1}%
  \noindent{\LARGE\bfseries\lblDia~#1}\quad
  {\large(\lblSemana~#2 \textperiodcentered\ \lblTrimestre~#3)}\par
  {\normalsize\bfseries\color{colorGris}#4}\par
  \vspace{1mm}
  \textbf{\lblFecha:} \rule{22mm}{0.4pt}\hspace{4mm}%
  \textbf{\lblHecho:} \casilla\ \lblHechoSola\quad
  \casilla\ \lblHechoConAyuda\quad
  \casilla\ \lblHechoConDificultad\hspace{4mm}%
  \textbf{\lblNotas:} \rule{20mm}{0.4pt}\par
  \vspace{3mm}
}

\makeatletter
\def\ps@pienumeros{%
  \let\@oddhead\@empty\let\@evenhead\@empty
  \def\@oddfoot{\hfil\footnotesize\color{colorGris}\lblDia~\diapagina@actual~/ \totaldias\hfil}%
  \let\@evenfoot\@oddfoot
}
\newenvironment{diapagina}[4]{%
  \gdef\diapagina@actual{#1}%
  \thispagestyle{pienumeros}%
  \cabeceraDia{#1}{#2}{#3}{#4}%
}{%
  \par\newpage
}

% Una etiqueta que solo guarda la página en la que cae, para
% tools/check_pages.py: la clave de respuestas empieza con
% \etiquetaPagina{clave}, así se comprueba también que el último día no
% se ha desbordado a una segunda página.
\newcommand{\etiquetaPagina}[1]{\begingroup\def\@currentlabel{}\label{#1}\endgroup}
\makeatother

% =========================================================================
% Hoy
% =========================================================================
% La caja de arriba: el dibujo del día (#1, un tikzpicture que compone
% tools/gen_economia.py) y la historia (#2), que lee la niña o el niño.
% Los lunes, debajo, la palabra de la semana (#3: \palabraSemana, o
% nada).
\newcommand{\hoy}[3]{%
  \begin{cajaHoy}
  \begin{tabularx}{\linewidth}{@{}>{\centering\arraybackslash}m{30mm}@{\hspace{5mm}}>{\sinPartir\arraybackslash\large}X@{}}
  #1 & #2 \\
  \end{tabularx}
  #3
  \end{cajaHoy}\par\vspace{4mm}}
% La palabra de la semana: la palabra (#1) y lo que quiere decir (#2).
\newtcolorbox{cajaPalabra}{enhanced, colback=white, colframe=colorLectura, boxrule=0.8pt,
  arc=2mm, left=3mm, right=3mm, top=1.5mm, bottom=1.5mm, before skip=2.5mm, after skip=0mm}
\newcommand{\palabraSemana}[2]{%
  \begin{cajaPalabra}\normalsize\textbf{\color{colorLectura}\lblPalabraSemana:}
  \textbf{\large #1}. #2\end{cajaPalabra}}

% =========================================================================
% Medalla de fin de trimestre (tools/gen_economia.py, PLANTILLA_MEDALLA)
% =========================================================================
\newcommand{\medalla}[1]{%
  \begin{tikzpicture}[line width=1.6pt, color=colorCrea]
    \draw (0,0) circle (1.1);
    \draw (0,0) circle (0.85);
    \node at (0,0) {\Large\bfseries #1};
    \draw (-0.35,-1.05) -- (-0.7,-2.2) -- (-0.15,-1.75) -- cycle;
    \draw (0.35,-1.05)  -- (0.7,-2.2)  -- (0.15,-1.75)  -- cycle;
  \end{tikzpicture}}

% =========================================================================
% La clave de respuestas y el diccionario
% =========================================================================
\newcommand{\claveEntrada}[3]{%
  \item \textbf{\lblDia\ #1} \textcolor{colorGris}{(#2)}: #3%
}
% Una palabra de "Mi diccionario de economía": la palabra (#1), lo que
% quiere decir (#2) y la semana en que salió (#3).
\newcommand{\entradaDiccionario}[3]{%
  \item \textbf{\color{colorLectura}#1}. #2 \textcolor{colorGris}{(\lblSemana~#3)}%
}
, y los mismos nombres de color
% apuntan a negro y gris: la misma página, sin tinta de color.
\makeatletter
\newif\ifbwbook
\def\aln@bookcolor@bw{bw}
\@ifundefined{bookcolor}{\bwbookfalse}{%
  \ifx\bookcolor\aln@bookcolor@bw \bwbooktrue \else \bwbookfalse \fi
}
\makeatother

\ifbwbook
  \definecolor{colorLectura}{gray}{0}
  \definecolor{colorLecturaClaro}{gray}{0.95}
  \definecolor{colorDibuja}{gray}{0}
  \definecolor{colorDibujaClaro}{gray}{0.95}
  \definecolor{colorCompleta}{gray}{0}
  \definecolor{colorCompletaClaro}{gray}{0.95}
  \definecolor{colorResponde}{gray}{0}
  \definecolor{colorRespondeClaro}{gray}{0.95}
  \definecolor{colorRelaciona}{gray}{0}
  \definecolor{colorRelacionaClaro}{gray}{0.95}
  \definecolor{colorCrea}{gray}{0}
  \definecolor{colorCreaClaro}{gray}{0.95}
  \definecolor{colorGris}{gray}{0.38}
  \definecolor{colorPunto}{gray}{0.25}
\else
  \definecolor{colorLectura}{HTML}{2E7D32}
  \definecolor{colorLecturaClaro}{HTML}{EAF6EC}
  \definecolor{colorDibuja}{HTML}{1565C0}
  \definecolor{colorDibujaClaro}{HTML}{E7F0FB}
  \definecolor{colorCompleta}{HTML}{C2185B}
  \definecolor{colorCompletaClaro}{HTML}{FBE9EF}
  \definecolor{colorResponde}{HTML}{6A1B9A}
  \definecolor{colorRespondeClaro}{HTML}{F2E9F7}
  \definecolor{colorRelaciona}{HTML}{E65100}
  \definecolor{colorRelacionaClaro}{HTML}{FDEEE2}
  \definecolor{colorCrea}{HTML}{AD1457}
  \definecolor{colorCreaClaro}{HTML}{FBE7EE}
  \definecolor{colorGris}{HTML}{616161}
  \definecolor{colorPunto}{HTML}{6A1B9A}
\fi

% =========================================================================
% Los dibujos
% =========================================================================
% diagrams/kit.tex: las piezas de los dibujos de «First Words» (de
% "Aprendo a leer"); diagrams/objetos.tex: las cosas de "Aprendo los
% números", cada una en la misma caja, con las monedas y los billetes;
% diagrams/economia.tex: las cosas nuevas de este cuaderno (la lupa, los
% cromos, las canicas, la hucha...).
\input{diagrams/kit}
\input{diagrams/objetos}
% diagrams/economia.tex -- las cosas nuevas de "Aprendo economía", en la
% misma caja que las de diagrams/objetos.tex (de -1 a 1 en x y en y),
% con el mismo estilo de línea: todas las zonas cerradas y rellenas de
% blanco, para que se puedan colorear.

% La lupa de la abuela Rosa (la de "Leo con lupa"): la lente, y el mango
% de madera con sus tres rayas grabadas.
\newcommand{\objLupa}{%
  \begin{scope}[rotate=45]
    \filldraw[fill=white, rounded corners=0.8mm] (-0.14,-1.12) rectangle (0.14,-0.52);
    \draw (-0.14,-0.72) -- (0.14,-0.72);
    \draw (-0.14,-0.84) -- (0.14,-0.84);
    \draw (-0.14,-0.96) -- (0.14,-0.96);
    \filldraw[fill=white] (-0.09,-0.54) rectangle (0.09,-0.42);
    \filldraw[fill=white] (0,0.12) circle (0.56);
    \draw (0,0.12) circle (0.45);
    \draw (-0.3,0.3) arc (145:195:0.33);
  \end{scope}}

% Un cromo: la tarjeta, el marco de la foto (una estrella) y dos líneas
% de texto.
\newcommand{\objCromo}{%
  \filldraw[fill=white, rounded corners=1mm] (-0.62,-0.9) rectangle (0.62,0.9);
  \draw[rounded corners=0.6mm] (-0.46,-0.3) rectangle (0.46,0.74);
  \filldraw[fill=white, line join=round, shift={(0,0.22)}]
    (90:0.34) -- (126:0.14) -- (162:0.34) -- (198:0.14) -- (234:0.34)
    -- (270:0.14) -- (306:0.34) -- (342:0.14) -- (18:0.34) -- (54:0.14) -- cycle;
  \draw (-0.42,-0.5) -- (0.42,-0.5);
  \draw (-0.42,-0.7) -- (0.18,-0.7);}

% Una canica, con su remolino y un brillo.
\newcommand{\objCanica}{%
  \filldraw[fill=white] (0,0) circle (0.72);
  \draw (-0.55,-0.18) .. controls (-0.2,0.35) and (0.22,-0.38) .. (0.58,0.2);
  \draw (-0.3,-0.55) .. controls (0,-0.25) and (0.2,-0.6) .. (0.45,-0.45);
  \draw (-0.42,0.28) arc (150:105:0.42);}

% Una hucha de cerdito: las patas, el cuerpo, el hocico, la oreja, el
% ojo, la ranura de las monedas y la cola.
\newcommand{\objHucha}{%
  \begin{scope}[shift={(-0.08,0.02)}]
    \filldraw[fill=white] (-0.5,-0.74) rectangle (-0.27,-0.42);
    \filldraw[fill=white] (0.28,-0.74) rectangle (0.51,-0.42);
    \filldraw[fill=white] (0,0) ellipse [x radius=0.84, y radius=0.6];
    \filldraw[fill=white] (0.86,-0.05) ellipse [x radius=0.15, y radius=0.21];
    \fill (0.83,-0.1) circle (0.03); \fill (0.9,0.02) circle (0.03);
    \filldraw[fill=white, line join=round] (0.3,0.48) -- (0.46,0.84) -- (0.6,0.4) -- cycle;
    \fill (0.52,0.16) circle (0.055);
    \draw[line width=1.8pt] (-0.28,0.52) -- (0.12,0.58);
    \draw (-0.84,0.04) .. controls (-1.0,0.14) and (-0.96,0.36) .. (-0.82,0.3);
  \end{scope}}


% =========================================================================
% Cajas
% =========================================================================
% La caja de actividad llena todo lo que queda de página (`height fill`),
% como en los otros cuadernos. Lo que tiene que ir abajo del todo (la
% línea de escribir, el cartel de cada diez páginas) va en la parte de
% abajo de la caja (\tcblower, estilo `abajo`); lo que tiene que ir
% centrado en el hueco, con `centrado abajo`.
\tcbset{
  cajabase/.style={
    enhanced,
    boxrule=1.1pt,
    arc=3mm,
    left=6mm, right=6mm, top=5mm, bottom=5mm,
    fonttitle=\bfseries\sffamily\large,
    coltitle=white,
  },
  cajadia/.style={cajabase, height fill},
  abajo/.style={space to upper, segmentation hidden, middle=2mm},
  centrado abajo/.style={space to lower, valign lower=center, segmentation hidden, middle=2mm},
}

\newtcolorbox{cajaHoy}{cajabase, top=3mm, bottom=3mm,
  colback=colorLecturaClaro, colframe=colorLectura, colbacktitle=colorLectura,
  title=\lblHoy}

% Una caja de actividad: #1 el nombre (cajaClasifica...), #2 el título
% (\lblClasifica...), #3 el color.
\newcommand{\cajaActividad}[3]{%
  \newtcolorbox{#1}[1][]{cajadia, ##1,
    colback=#3Claro, colframe=#3, colbacktitle=#3, title=#2}}
\cajaActividad{cajaClasifica}{\lblClasifica}{colorRelaciona}
\cajaActividad{cajaUne}{\lblUne}{colorRelaciona}
\cajaActividad{cajaVF}{\lblVF}{colorResponde}
\cajaActividad{cajaTrueque}{\lblTrueque}{colorCompleta}
\cajaActividad{cajaDibuja}{\lblDibuja}{colorDibuja}
\cajaActividad{cajaRepasa}{\lblRepasa}{colorCrea}

% La instrucción de cada actividad: la lee el adulto.
\newcommand{\instruccion}[1]{{\footnotesize\color{colorGris}#1\par}}
% El enunciado grande, lo que se hace: "Une cada cosa con su caja." Sin
% justificar y sin partir palabras, como la historia de "Hoy": lo lee la
% niña o el niño.
\newcommand{\sinCortes}{\hyphenpenalty=10000\exhyphenpenalty=10000}
\newcommand{\sinPartir}{\raggedright\sinCortes}
\newcommand{\enunciado}[1]{{\LARGE\sinPartir #1\par}}
% El hueco en el que se escribe la respuesta (#1, el ancho).
\newcommand{\huecoRespuesta}[1][26mm]{\raisebox{-5mm}{\tikz\draw[line width=1.1pt,
  rounded corners=2mm, fill=white] (0,0) rectangle (#1,18mm);}}
% Lo que no cabe en su sitio, encogido hasta que quepa (#1, el ancho).
\newsavebox{\cajaAjustada}
\newcommand{\ajustado}[2]{\sbox{\cajaAjustada}{#2}%
  \ifdim\wd\cajaAjustada>#1\relax\resizebox{#1}{!}{\usebox{\cajaAjustada}}%
  \else\usebox{\cajaAjustada}\fi}
% El cartel de cada diez páginas, al pie de "Repasa".
\newcommand{\cartel}[1]{\begin{center}\bfseries\Large\color{colorCrea}#1\end{center}}

\newcommand{\casilla}{\raisebox{-0.15ex}{\framebox[3.4mm][c]{\rule{0pt}{2.6mm}}}}
\newenvironment{listaRepaso}{%
  \begin{itemize}[label=\casilla, leftmargin=8mm, itemsep=2mm, topsep=1mm]\large
}{%
  \end{itemize}
}

% "¿Verdad o mentira?": cada frase, con su V y su F para rodear.
\newcommand{\letraVF}[1]{\tikz[baseline=-1.1mm]{\draw[line width=1pt] (0,0) circle (4.2mm);
  \node[font=\Large\bfseries] at (0,0) {#1};}}
\newcommand{\frasevf}[2]{%
  \noindent\begin{tabularx}{\linewidth}{@{}r@{\hspace{3mm}}>{\sinPartir\arraybackslash\Large}X@{\hspace{5mm}}c@{\hspace{3mm}}c@{}}
  {\Large\bfseries #1.} & #2 & \letraVF{\lblVerdad} & \letraVF{\lblMentira} \\
  \end{tabularx}\par\vspace{9mm}}

% =========================================================================
% La página de un día
% =========================================================================
% La cabecera es la de "Aprendo a leer": el día, la semana y el
% trimestre; el tema de la semana (el mismo que esa semana en "Leo con
% lupa"), y los campos para el adulto. Deja \label{dia:N}:
% tools/check_pages.py comprueba con esas etiquetas que cada día ocupa
% exactamente una página.
\newcommand{\cabeceraDia}[4]{%
  \label{dia:#1}%
  \noindent{\LARGE\bfseries\lblDia~#1}\quad
  {\large(\lblSemana~#2 \textperiodcentered\ \lblTrimestre~#3)}\par
  {\normalsize\bfseries\color{colorGris}#4}\par
  \vspace{1mm}
  \textbf{\lblFecha:} \rule{22mm}{0.4pt}\hspace{4mm}%
  \textbf{\lblHecho:} \casilla\ \lblHechoSola\quad
  \casilla\ \lblHechoConAyuda\quad
  \casilla\ \lblHechoConDificultad\hspace{4mm}%
  \textbf{\lblNotas:} \rule{20mm}{0.4pt}\par
  \vspace{3mm}
}

\makeatletter
\def\ps@pienumeros{%
  \let\@oddhead\@empty\let\@evenhead\@empty
  \def\@oddfoot{\hfil\footnotesize\color{colorGris}\lblDia~\diapagina@actual~/ \totaldias\hfil}%
  \let\@evenfoot\@oddfoot
}
\newenvironment{diapagina}[4]{%
  \gdef\diapagina@actual{#1}%
  \thispagestyle{pienumeros}%
  \cabeceraDia{#1}{#2}{#3}{#4}%
}{%
  \par\newpage
}

% Una etiqueta que solo guarda la página en la que cae, para
% tools/check_pages.py: la clave de respuestas empieza con
% \etiquetaPagina{clave}, así se comprueba también que el último día no
% se ha desbordado a una segunda página.
\newcommand{\etiquetaPagina}[1]{\begingroup\def\@currentlabel{}\label{#1}\endgroup}
\makeatother

% =========================================================================
% Hoy
% =========================================================================
% La caja de arriba: el dibujo del día (#1, un tikzpicture que compone
% tools/gen_economia.py) y la historia (#2), que lee la niña o el niño.
% Los lunes, debajo, la palabra de la semana (#3: \palabraSemana, o
% nada).
\newcommand{\hoy}[3]{%
  \begin{cajaHoy}
  \begin{tabularx}{\linewidth}{@{}>{\centering\arraybackslash}m{30mm}@{\hspace{5mm}}>{\sinPartir\arraybackslash\large}X@{}}
  #1 & #2 \\
  \end{tabularx}
  #3
  \end{cajaHoy}\par\vspace{4mm}}
% La palabra de la semana: la palabra (#1) y lo que quiere decir (#2).
\newtcolorbox{cajaPalabra}{enhanced, colback=white, colframe=colorLectura, boxrule=0.8pt,
  arc=2mm, left=3mm, right=3mm, top=1.5mm, bottom=1.5mm, before skip=2.5mm, after skip=0mm}
\newcommand{\palabraSemana}[2]{%
  \begin{cajaPalabra}\normalsize\textbf{\color{colorLectura}\lblPalabraSemana:}
  \textbf{\large #1}. #2\end{cajaPalabra}}

% =========================================================================
% Medalla de fin de trimestre (tools/gen_economia.py, PLANTILLA_MEDALLA)
% =========================================================================
\newcommand{\medalla}[1]{%
  \begin{tikzpicture}[line width=1.6pt, color=colorCrea]
    \draw (0,0) circle (1.1);
    \draw (0,0) circle (0.85);
    \node at (0,0) {\Large\bfseries #1};
    \draw (-0.35,-1.05) -- (-0.7,-2.2) -- (-0.15,-1.75) -- cycle;
    \draw (0.35,-1.05)  -- (0.7,-2.2)  -- (0.15,-1.75)  -- cycle;
  \end{tikzpicture}}

% =========================================================================
% La clave de respuestas y el diccionario
% =========================================================================
\newcommand{\claveEntrada}[3]{%
  \item \textbf{\lblDia\ #1} \textcolor{colorGris}{(#2)}: #3%
}
% Una palabra de "Mi diccionario de economía": la palabra (#1), lo que
% quiere decir (#2) y la semana en que salió (#3).
\newcommand{\entradaDiccionario}[3]{%
  \item \textbf{\color{colorLectura}#1}. #2 \textcolor{colorGris}{(\lblSemana~#3)}%
}
, y los mismos nombres de color
% apuntan a negro y gris: la misma página, sin tinta de color.
\makeatletter
\newif\ifbwbook
\def\aln@bookcolor@bw{bw}
\@ifundefined{bookcolor}{\bwbookfalse}{%
  \ifx\bookcolor\aln@bookcolor@bw \bwbooktrue \else \bwbookfalse \fi
}
\makeatother

\ifbwbook
  \definecolor{colorLectura}{gray}{0}
  \definecolor{colorLecturaClaro}{gray}{0.95}
  \definecolor{colorDibuja}{gray}{0}
  \definecolor{colorDibujaClaro}{gray}{0.95}
  \definecolor{colorCompleta}{gray}{0}
  \definecolor{colorCompletaClaro}{gray}{0.95}
  \definecolor{colorResponde}{gray}{0}
  \definecolor{colorRespondeClaro}{gray}{0.95}
  \definecolor{colorRelaciona}{gray}{0}
  \definecolor{colorRelacionaClaro}{gray}{0.95}
  \definecolor{colorCrea}{gray}{0}
  \definecolor{colorCreaClaro}{gray}{0.95}
  \definecolor{colorGris}{gray}{0.38}
  \definecolor{colorPunto}{gray}{0.25}
\else
  \definecolor{colorLectura}{HTML}{2E7D32}
  \definecolor{colorLecturaClaro}{HTML}{EAF6EC}
  \definecolor{colorDibuja}{HTML}{1565C0}
  \definecolor{colorDibujaClaro}{HTML}{E7F0FB}
  \definecolor{colorCompleta}{HTML}{C2185B}
  \definecolor{colorCompletaClaro}{HTML}{FBE9EF}
  \definecolor{colorResponde}{HTML}{6A1B9A}
  \definecolor{colorRespondeClaro}{HTML}{F2E9F7}
  \definecolor{colorRelaciona}{HTML}{E65100}
  \definecolor{colorRelacionaClaro}{HTML}{FDEEE2}
  \definecolor{colorCrea}{HTML}{AD1457}
  \definecolor{colorCreaClaro}{HTML}{FBE7EE}
  \definecolor{colorGris}{HTML}{616161}
  \definecolor{colorPunto}{HTML}{6A1B9A}
\fi

% =========================================================================
% Los dibujos
% =========================================================================
% diagrams/kit.tex: las piezas de los dibujos de «First Words» (de
% "Aprendo a leer"); diagrams/objetos.tex: las cosas de "Aprendo los
% números", cada una en la misma caja, con las monedas y los billetes;
% diagrams/economia.tex: las cosas nuevas de este cuaderno (la lupa, los
% cromos, las canicas, la hucha...).
\input{diagrams/kit}
\input{diagrams/objetos}
% diagrams/economia.tex -- las cosas nuevas de "Aprendo economía", en la
% misma caja que las de diagrams/objetos.tex (de -1 a 1 en x y en y),
% con el mismo estilo de línea: todas las zonas cerradas y rellenas de
% blanco, para que se puedan colorear.

% La lupa de la abuela Rosa (la de "Leo con lupa"): la lente, y el mango
% de madera con sus tres rayas grabadas.
\newcommand{\objLupa}{%
  \begin{scope}[rotate=45]
    \filldraw[fill=white, rounded corners=0.8mm] (-0.14,-1.12) rectangle (0.14,-0.52);
    \draw (-0.14,-0.72) -- (0.14,-0.72);
    \draw (-0.14,-0.84) -- (0.14,-0.84);
    \draw (-0.14,-0.96) -- (0.14,-0.96);
    \filldraw[fill=white] (-0.09,-0.54) rectangle (0.09,-0.42);
    \filldraw[fill=white] (0,0.12) circle (0.56);
    \draw (0,0.12) circle (0.45);
    \draw (-0.3,0.3) arc (145:195:0.33);
  \end{scope}}

% Un cromo: la tarjeta, el marco de la foto (una estrella) y dos líneas
% de texto.
\newcommand{\objCromo}{%
  \filldraw[fill=white, rounded corners=1mm] (-0.62,-0.9) rectangle (0.62,0.9);
  \draw[rounded corners=0.6mm] (-0.46,-0.3) rectangle (0.46,0.74);
  \filldraw[fill=white, line join=round, shift={(0,0.22)}]
    (90:0.34) -- (126:0.14) -- (162:0.34) -- (198:0.14) -- (234:0.34)
    -- (270:0.14) -- (306:0.34) -- (342:0.14) -- (18:0.34) -- (54:0.14) -- cycle;
  \draw (-0.42,-0.5) -- (0.42,-0.5);
  \draw (-0.42,-0.7) -- (0.18,-0.7);}

% Una canica, con su remolino y un brillo.
\newcommand{\objCanica}{%
  \filldraw[fill=white] (0,0) circle (0.72);
  \draw (-0.55,-0.18) .. controls (-0.2,0.35) and (0.22,-0.38) .. (0.58,0.2);
  \draw (-0.3,-0.55) .. controls (0,-0.25) and (0.2,-0.6) .. (0.45,-0.45);
  \draw (-0.42,0.28) arc (150:105:0.42);}

% Una hucha de cerdito: las patas, el cuerpo, el hocico, la oreja, el
% ojo, la ranura de las monedas y la cola.
\newcommand{\objHucha}{%
  \begin{scope}[shift={(-0.08,0.02)}]
    \filldraw[fill=white] (-0.5,-0.74) rectangle (-0.27,-0.42);
    \filldraw[fill=white] (0.28,-0.74) rectangle (0.51,-0.42);
    \filldraw[fill=white] (0,0) ellipse [x radius=0.84, y radius=0.6];
    \filldraw[fill=white] (0.86,-0.05) ellipse [x radius=0.15, y radius=0.21];
    \fill (0.83,-0.1) circle (0.03); \fill (0.9,0.02) circle (0.03);
    \filldraw[fill=white, line join=round] (0.3,0.48) -- (0.46,0.84) -- (0.6,0.4) -- cycle;
    \fill (0.52,0.16) circle (0.055);
    \draw[line width=1.8pt] (-0.28,0.52) -- (0.12,0.58);
    \draw (-0.84,0.04) .. controls (-1.0,0.14) and (-0.96,0.36) .. (-0.82,0.3);
  \end{scope}}


% =========================================================================
% Cajas
% =========================================================================
% La caja de actividad llena todo lo que queda de página (`height fill`),
% como en los otros cuadernos. Lo que tiene que ir abajo del todo (la
% línea de escribir, el cartel de cada diez páginas) va en la parte de
% abajo de la caja (\tcblower, estilo `abajo`); lo que tiene que ir
% centrado en el hueco, con `centrado abajo`.
\tcbset{
  cajabase/.style={
    enhanced,
    boxrule=1.1pt,
    arc=3mm,
    left=6mm, right=6mm, top=5mm, bottom=5mm,
    fonttitle=\bfseries\sffamily\large,
    coltitle=white,
  },
  cajadia/.style={cajabase, height fill},
  abajo/.style={space to upper, segmentation hidden, middle=2mm},
  centrado abajo/.style={space to lower, valign lower=center, segmentation hidden, middle=2mm},
}

\newtcolorbox{cajaHoy}{cajabase, top=3mm, bottom=3mm,
  colback=colorLecturaClaro, colframe=colorLectura, colbacktitle=colorLectura,
  title=\lblHoy}

% Una caja de actividad: #1 el nombre (cajaClasifica...), #2 el título
% (\lblClasifica...), #3 el color.
\newcommand{\cajaActividad}[3]{%
  \newtcolorbox{#1}[1][]{cajadia, ##1,
    colback=#3Claro, colframe=#3, colbacktitle=#3, title=#2}}
\cajaActividad{cajaClasifica}{\lblClasifica}{colorRelaciona}
\cajaActividad{cajaUne}{\lblUne}{colorRelaciona}
\cajaActividad{cajaVF}{\lblVF}{colorResponde}
\cajaActividad{cajaTrueque}{\lblTrueque}{colorCompleta}
\cajaActividad{cajaDibuja}{\lblDibuja}{colorDibuja}
\cajaActividad{cajaRepasa}{\lblRepasa}{colorCrea}

% La instrucción de cada actividad: la lee el adulto.
\newcommand{\instruccion}[1]{{\footnotesize\color{colorGris}#1\par}}
% El enunciado grande, lo que se hace: "Une cada cosa con su caja." Sin
% justificar y sin partir palabras, como la historia de "Hoy": lo lee la
% niña o el niño.
\newcommand{\sinCortes}{\hyphenpenalty=10000\exhyphenpenalty=10000}
\newcommand{\sinPartir}{\raggedright\sinCortes}
\newcommand{\enunciado}[1]{{\LARGE\sinPartir #1\par}}
% El hueco en el que se escribe la respuesta (#1, el ancho).
\newcommand{\huecoRespuesta}[1][26mm]{\raisebox{-5mm}{\tikz\draw[line width=1.1pt,
  rounded corners=2mm, fill=white] (0,0) rectangle (#1,18mm);}}
% Lo que no cabe en su sitio, encogido hasta que quepa (#1, el ancho).
\newsavebox{\cajaAjustada}
\newcommand{\ajustado}[2]{\sbox{\cajaAjustada}{#2}%
  \ifdim\wd\cajaAjustada>#1\relax\resizebox{#1}{!}{\usebox{\cajaAjustada}}%
  \else\usebox{\cajaAjustada}\fi}
% El cartel de cada diez páginas, al pie de "Repasa".
\newcommand{\cartel}[1]{\begin{center}\bfseries\Large\color{colorCrea}#1\end{center}}

\newcommand{\casilla}{\raisebox{-0.15ex}{\framebox[3.4mm][c]{\rule{0pt}{2.6mm}}}}
\newenvironment{listaRepaso}{%
  \begin{itemize}[label=\casilla, leftmargin=8mm, itemsep=2mm, topsep=1mm]\large
}{%
  \end{itemize}
}

% "¿Verdad o mentira?": cada frase, con su V y su F para rodear.
\newcommand{\letraVF}[1]{\tikz[baseline=-1.1mm]{\draw[line width=1pt] (0,0) circle (4.2mm);
  \node[font=\Large\bfseries] at (0,0) {#1};}}
\newcommand{\frasevf}[2]{%
  \noindent\begin{tabularx}{\linewidth}{@{}r@{\hspace{3mm}}>{\sinPartir\arraybackslash\Large}X@{\hspace{5mm}}c@{\hspace{3mm}}c@{}}
  {\Large\bfseries #1.} & #2 & \letraVF{\lblVerdad} & \letraVF{\lblMentira} \\
  \end{tabularx}\par\vspace{9mm}}

% =========================================================================
% La página de un día
% =========================================================================
% La cabecera es la de "Aprendo a leer": el día, la semana y el
% trimestre; el tema de la semana (el mismo que esa semana en "Leo con
% lupa"), y los campos para el adulto. Deja \label{dia:N}:
% tools/check_pages.py comprueba con esas etiquetas que cada día ocupa
% exactamente una página.
\newcommand{\cabeceraDia}[4]{%
  \label{dia:#1}%
  \noindent{\LARGE\bfseries\lblDia~#1}\quad
  {\large(\lblSemana~#2 \textperiodcentered\ \lblTrimestre~#3)}\par
  {\normalsize\bfseries\color{colorGris}#4}\par
  \vspace{1mm}
  \textbf{\lblFecha:} \rule{22mm}{0.4pt}\hspace{4mm}%
  \textbf{\lblHecho:} \casilla\ \lblHechoSola\quad
  \casilla\ \lblHechoConAyuda\quad
  \casilla\ \lblHechoConDificultad\hspace{4mm}%
  \textbf{\lblNotas:} \rule{20mm}{0.4pt}\par
  \vspace{3mm}
}

\makeatletter
\def\ps@pienumeros{%
  \let\@oddhead\@empty\let\@evenhead\@empty
  \def\@oddfoot{\hfil\footnotesize\color{colorGris}\lblDia~\diapagina@actual~/ \totaldias\hfil}%
  \let\@evenfoot\@oddfoot
}
\newenvironment{diapagina}[4]{%
  \gdef\diapagina@actual{#1}%
  \thispagestyle{pienumeros}%
  \cabeceraDia{#1}{#2}{#3}{#4}%
}{%
  \par\newpage
}

% Una etiqueta que solo guarda la página en la que cae, para
% tools/check_pages.py: la clave de respuestas empieza con
% \etiquetaPagina{clave}, así se comprueba también que el último día no
% se ha desbordado a una segunda página.
\newcommand{\etiquetaPagina}[1]{\begingroup\def\@currentlabel{}\label{#1}\endgroup}
\makeatother

% =========================================================================
% Hoy
% =========================================================================
% La caja de arriba: el dibujo del día (#1, un tikzpicture que compone
% tools/gen_economia.py) y la historia (#2), que lee la niña o el niño.
% Los lunes, debajo, la palabra de la semana (#3: \palabraSemana, o
% nada).
\newcommand{\hoy}[3]{%
  \begin{cajaHoy}
  \begin{tabularx}{\linewidth}{@{}>{\centering\arraybackslash}m{30mm}@{\hspace{5mm}}>{\sinPartir\arraybackslash\large}X@{}}
  #1 & #2 \\
  \end{tabularx}
  #3
  \end{cajaHoy}\par\vspace{4mm}}
% La palabra de la semana: la palabra (#1) y lo que quiere decir (#2).
\newtcolorbox{cajaPalabra}{enhanced, colback=white, colframe=colorLectura, boxrule=0.8pt,
  arc=2mm, left=3mm, right=3mm, top=1.5mm, bottom=1.5mm, before skip=2.5mm, after skip=0mm}
\newcommand{\palabraSemana}[2]{%
  \begin{cajaPalabra}\normalsize\textbf{\color{colorLectura}\lblPalabraSemana:}
  \textbf{\large #1}. #2\end{cajaPalabra}}

% =========================================================================
% Medalla de fin de trimestre (tools/gen_economia.py, PLANTILLA_MEDALLA)
% =========================================================================
\newcommand{\medalla}[1]{%
  \begin{tikzpicture}[line width=1.6pt, color=colorCrea]
    \draw (0,0) circle (1.1);
    \draw (0,0) circle (0.85);
    \node at (0,0) {\Large\bfseries #1};
    \draw (-0.35,-1.05) -- (-0.7,-2.2) -- (-0.15,-1.75) -- cycle;
    \draw (0.35,-1.05)  -- (0.7,-2.2)  -- (0.15,-1.75)  -- cycle;
  \end{tikzpicture}}

% =========================================================================
% La clave de respuestas y el diccionario
% =========================================================================
\newcommand{\claveEntrada}[3]{%
  \item \textbf{\lblDia\ #1} \textcolor{colorGris}{(#2)}: #3%
}
% Una palabra de "Mi diccionario de economía": la palabra (#1), lo que
% quiere decir (#2) y la semana en que salió (#3).
\newcommand{\entradaDiccionario}[3]{%
  \item \textbf{\color{colorLectura}#1}. #2 \textcolor{colorGris}{(\lblSemana~#3)}%
}
, y preamble.tex lee este fichero en vez de
% lang/es.tex. Lo que cambia cada día (los enunciados, las instrucciones
% y la clave) lo escribe tools/gen_economia.py con tools/idiomas.py.
% -----------------------------------------------------------------------

% Total de días del cuaderno. tools/gen_economia.py --check comprueba que
% coincide con su propio TOTAL_DIAS.
\newcommand{\totaldias}{260}

% La cabecera de cada día (\cabeceraDia, preamble.tex). "Część", como
% "Part" en inglés: cada parte es una estación.
\newcommand{\lblFecha}{Data}
\newcommand{\lblNotas}{Uwagi}
\newcommand{\lblDia}{Dzień}
\newcommand{\lblSemana}{Tydzień}
\newcommand{\lblTrimestre}{Część}
\newcommand{\lblHecho}{Zrobione}
\newcommand{\lblHechoSola}{samodzielnie}
\newcommand{\lblHechoConAyuda}{z pomocą}
\newcommand{\lblHechoConDificultad}{z trudem}

% La caja de arriba, la misma todos los días.
\newcommand{\lblHoy}{Dzisiaj}
\newcommand{\lblPalabraSemana}{Słowo tygodnia}

% Los títulos de las cajas de actividad.
\newcommand{\lblClasifica}{Posortuj}
\newcommand{\lblUne}{Połącz}
\newcommand{\lblVF}{Prawda czy fałsz?}
\newcommand{\lblTrueque}{Wymiana}
\newcommand{\lblDibuja}{Narysuj}
\newcommand{\lblRepasa}{Powtórka}
\newcommand{\lblDinero}{Pieniądze}
\newcommand{\lblReparte}{Podziel}
\newcommand{\lblCompra}{Zakupy}
\newcommand{\lblProblema}{Zadanie}
\newcommand{\lblBotes}{Trzy słoiki}
\newcommand{\lblHucha}{Skarbonka}
\newcommand{\lblLlega}{Czy wystarczy?}
\newcommand{\lblVuelta}{Reszta}
\newcommand{\lblOrdena}{Ułóż po kolei}
\newcommand{\lblElige}{Wybierz}
\newcommand{\lblCompara}{Porównaj}
\newcommand{\lblCuentas}{Moje rachunki}
% Las columnas del cuaderno de cuentas: las del dinero, de 24 mm, y en
% polaco las palabras son más largas (\ajustado, en preamble.tex, las
% encoge si no caben).
\newcommand{\lblQuePasa}{Co się dzieje}
\newcommand{\lblEntra}{\ajustado{23mm}{Dostaję}}
\newcommand{\lblSale}{\ajustado{23mm}{Wydaję}}
\newcommand{\lblQueda}{Zostaje}
\newcommand{\lblAlEmpezar}{Na początku}

% El dinero en los dibujos (las etiquetas de precio, las monedas y los
% billetes): el número y, detrás, el euro, como en español.
\newcommand{\importe}[1]{#1\,€}

% Los nombres de los tres botes (el tarro mide 30 mm).
\newcommand{\lblAhorrar}{\ajustado{27mm}{Oszczędzać}}
\newcommand{\lblGastar}{\ajustado{27mm}{Wydawać}}
\newcommand{\lblCompartir}{\ajustado{27mm}{Dzielić się}}

% Las dos letras de "Prawda czy fałsz?".
\newcommand{\lblVerdad}{P}
\newcommand{\lblMentira}{F}

% La portada.
\newcommand{\lblTituloPortada}{Poznaję ekonomię}
\newcommand{\lblSubtituloPortada}{Pieniądze, praca i codzienne wybory}
\newcommand{\lblSubtituloPortadaMetodo}{Zeszyt na każdy dzień: uczę się oszczędzać, wybierać, wymieniać się i dzielić -- razem z Lucíą i jej rodziną}
\newcommand{\lblPortadaNombre}{Imię}
\newcommand{\lblPortadaCurso}{Klasa}

% La medalla del final de cada parte (T1-T3) y el cartel de cada diez
% páginas ("Już 10 stron!": con 10, 20, 30..., siempre "stron").
\newcommand{\lblMedalla}[1]{Medal za #1!}
\newcommand{\lblMedallaAnimo}{Tak trzymaj, \rule{55mm}{0.4pt}!}
\newcommand{\lblPaginasHechas}[1]{Już #1 stron!}

% La clave de respuestas y el diccionario.
\newcommand{\lblClave}{Odpowiedzi}
\newcommand{\lblDiccionario}{Mój słowniczek ekonomiczny}
